По окончании высшего учебного заведения каждый выпускник сталкивается с определением траектории своего дальнейшего профессионального развития. В связи с этим индивидуальное домашнее задание по дисциплине «Тайм-менеджмент» выполняется на тему «Разработка плана профессионального роста».

Среди очевидных вариантов в настоящее время можно выделить занятость по найму на условиях Трудового кодекса Российской Федерации и работа во фрилансе с оформлением в качестве плательщика налога на профессиональный доход (самозанятость). 

В связи с этим, в индивидуальном домашнем задании обучающимся предлагается оценить пути возможного трудоустройства, 

Задание включает разработку трех разделов. 

Первый раздел предполагает изучение своих личностных особенностей: текущих знаний и навыков, интересов и склонностей, свойств характера. Результатом самоанализа становится формулировка двух альтернативных путей профессионального развития.

Второй раздел включает в себя анализ первого альтернативного пути профессионального развития – трудоустройства в компанию. Раздел включает изучение требований рынка труда на основании опубликованных вакансий и профессионального стандарта, выявление разрыва и составление индивидуального плана развития по его устранению, анализ потенциального дохода, выбор компании для трудоустройства и составление резюме. 

Третий раздел включает в себя анализ второго альтернативного пути профессионального развития. В случае если он предполагает трудоустройство в компанию, то его содержание аналогично второму разделу, но для другой должности. Если же он предполагает самозанятость, то необходимо описать  своего будущего клиента и выстроить стратегию продвижения, проанализировать рынок фриланса по аналогичным услугам/продукции и рассчитать потенциальный доход.

% Четвертый раздел посвящен планированию реализации выбранной траектории развития посредством построения графика Ганта. В данном разделе обучающемуся необходимо систематизировать мероприятия индивидуального плана развития, определив их последовательность и продолжительность во времени. Построение графика требует учета текущей академической нагрузки, профессиональной деятельности (при наличии), а также иных значимых жизненных ролей для обеспечения сбалансированного распределения временных ресурсов

В заключении необходимо провести сравнение оцениваемых вариантов будущего трудоустройства, определить их преимущества и недостатки с учетом личных предпочтений, нагрузки, а также потенциального дохода.
